\documentclass[12pt]{article}

% Layout
\usepackage[margin=1.1in]{geometry}
\usepackage{setspace}
\setstretch{1.15}

% Math
\usepackage{amsmath, amssymb}
\usepackage{bm}

% Figures & Tables
\usepackage{graphicx}
\usepackage{float}
\usepackage{booktabs}
\usepackage{tabularx}
\usepackage{xcolor}

% Lists
\usepackage{enumitem}

% Typography & misc
\usepackage{microtype}
\usepackage{hyperref}
\hypersetup{colorlinks=true, linkcolor=blue!50!black, urlcolor=blue!50!black, citecolor=blue!50!black}

\title{Next Steps --- GCR Dosimetry Follow-On}
\author{Aryan Hemendra Shah}

\begin{document}
\maketitle
\textit{\small Last updated: \today}

\section{9/25/26 --- where things stand}

\texttt{gcrisk} (formerly \texttt{gcr-dosimetry-pipeline}) is submitted to
\textit{Life Sciences in Space Research}. Manuscript number assigned,
build-for-approval PDF came back clean --- no compile log dump, all 7
figures rendering with real data, all cross-refs/citations resolved. Also
found and fixed a real repo problem today: local \texttt{main} was 11
commits ahead of GitHub and 2 behind (a stale README edit made directly on
GitHub back on 6/4/26, before the \texttt{gcrisk} rename and the ISS/RAD
validation work). Merged, caught the merge silently reverting one sentence
to the old typo'd version, fixed it, pushed. Repo is clean and in sync now,
including the full EM submission package (\texttt{paper/submission/}) so
the exact files submitted are preserved in version control.

So: nothing blocking, no open repo issues. This document is just about
what's worth doing next, not anything urgent.

\section{9/25/26 --- skimmed two ScienceDirect exports for next-project ideas}

Two unrelated clusters, both from the same \textit{LSSR} special issue:

\begin{enumerate}[itemsep=3pt]
    \item \textbf{Astrochemistry / prebiotic ice chemistry} (13 papers) ---
    COMs in protostellar systems, nucleobases in interstellar ices, glycine
    stability under ionizing radiation, cosmic-ray implantation into
    molecular clouds. Interesting field, but a genuinely different research
    track from anything \texttt{gcrisk} does. Not pursuing unless I
    deliberately want to branch into astrobiology/astrochemistry, which is
    a separate decision from ``what's the sequel to the GCR paper.''

    \item \textbf{Space radiation dose-response / quality-factor modeling}
    (14 papers) --- this is the one that matters. Several of these are
    direct hits against the exact uncertainty structure the just-submitted
    paper already quantified.
\end{enumerate}

The overlap that matters: the paper's own value-of-information sweep found
that halving $Q_{\text{factor}}$ uncertainty alone shrinks p95 REID by
33.5\% --- about 7$\times$ more leverage than any other single parameter,
more than every GCR transport-physics parameter combined. That's not a
throwaway finding, it's a pointer at exactly what to work on next, and this
paper batch landed with several candidate answers to it.

\section{Candidate next projects, ranked}

\begin{center}
\begin{tabularx}{\textwidth}{
  >{\raggedright\arraybackslash}p{3.1cm}
  >{\raggedright\arraybackslash}X
  >{\raggedright\arraybackslash}p{1.9cm}
}
\toprule
\textbf{Idea} & \textbf{What it is} & \textbf{Effort} \\
\midrule
Charge-weighted $Q^*$ &
Naito \& Kodaira propose $Q^*$, a GCR-specific quality factor based on LET
and charge that diverges from ICRP-60 $Q(L)$ at high LET, claiming current
dose assessments over-weight the heavy-ion contribution. Direct drop-in
alternative to what \texttt{reid.py} already uses. &
Low \\[6pt]
Ensemble dose-response model &
Slaba \& Poignant (NASA Langley) --- a probabilistic, non-linear QF
framework meant to replace the Cucinotta-era calibration approach the
pipeline currently inherits. This is the actual research-frontier answer,
not a swap-in. &
High \\[6pt]
ICRP $Q(L)$ history/meta-analysis &
Sato's review --- not a project, but a citation that directly backs why
$Q$-factor is the dominant uncertainty source. Should already be in the
lit review. &
Trivial \\[6pt]
Murine organ DCFs + gastric cancer &
Hosseini et al.\ (organ dose conversion factors for GCRsim mice) and Suman
et al.\ (dose/LET-dependent gastric tumor incidence in
Apc1638N/+ mice) --- real GCRsim-derived organ-dose and tumor-incidence
data. Interesting as a validation target if REID validation ever moves
past human epidemiological life tables into direct dose-to-tumor-incidence
data. &
Medium, bigger scope \\
\bottomrule
\end{tabularx}
\end{center}

\textbf{Pick: $Q^*$.} Best ratio of effort to payoff --- it's a direct,
quantifiable extension of the exact sensitivity result the submitted paper
already highlights, reuses the existing LHS/REID machinery almost
unchanged, and gives a natural ``part 2'' paper: same validated pipeline,
one swapped component, a concrete number for how much REID moves.

\section{Build on top of \texttt{gcrisk}, or start separate?}

\textbf{On top of \texttt{gcrisk}.} The whole value of the $Q^*$ idea is
that it's the same validated dose/REID/LHS pipeline with one component
swapped --- forking a new repo means re-deriving the transport and
validation harness for no benefit.

Risk: bolting a second quality-factor model directly into \texttt{reid.py}
muddies the story of a package the just-submitted paper already describes
in print. Mitigation: keep it in a clearly separate module rather than a
rewrite in place ---

\begin{itemize}[itemsep=2pt]
    \item \texttt{gcrisk/qf\_models.py} --- \texttt{qf\_icrp60(LET)} (existing
    logic, moved/kept as-is) and \texttt{qf\_naito\_kodaira(LET, Z)} (new).
    \item \texttt{reid.py} takes a \texttt{qf\_model} argument rather than
    hard-coding ICRP-60.
    \item Version bump/tag so the new paper can just say ``\texttt{gcrisk}
    vX.Y, run with QF model $=Q^*$'' --- core package stays legible, second
    paper stays citable and reproducible off a pinned version.
\end{itemize}

\section{Concrete next action}

Not started yet --- this is the actual to-do, next time I sit down with
this:

\begin{enumerate}[itemsep=3pt]
    \item Pull the exact $Q^*(\text{LET}, Z)$ formula and $Z$-contribution
    weighting out of the Naito \& Kodaira PDF (\textit{Charge-weighted
    quality factors based on LET for practical dose assessment of space
    radiation}).
    \item Implement \texttt{qf\_naito\_kodaira()} in a new
    \texttt{gcrisk/qf\_models.py}, unit-tested against the paper's own
    reported $Q^*$ values at a few LET points before touching the ensemble.
    \item Add the \texttt{qf\_model} switch to \texttt{reid.py} /
    \texttt{uncertainty.py} without changing default behavior (ICRP-60
    stays default).
    \item Re-run the N=500 LHS ensemble with $Q^*$ in place of ICRP-60 at
    the same 16\,g/cm$^2$ Al calibration geometry already used in the
    submitted paper.
    \item Report the shift: median REID, p5--p95 REID, and how much of the
    75\% $Q_{\text{factor}}$-driven variance share moves. That number is
    the whole paper.
\end{enumerate}

\end{document}
